El teorema de Cantor-Berstein asegura que si $X \sim Y_1 \subset Y$y $Y \sim X_1 \subset X$, entonces $X \sim Y$. \newline

\begin{itemize}
  \item Dem: Basta probar que existe un funcion \textbf{biyectiva}: Una funcion $f$se dice biyeccion si es inyectiva y suprayectiva. $h:X \to Y$. Para esto podemos usar las propiedades de las funciones     \textbf{inyectivas}: Una funcion $f$se dice inyectiva si: \newline
  $f(x)=f(y) \; \Rightarrow \; x=y$. Para ello, definimos las funciones $f: X \to Y$y $g: Y \to X$, ambas funciones inyectivas.      Antes de definir la funcion $h$, debemos definir dos conjuntos fundamentales, estos son $X_\infty$, este conjunto contiene todos los elementos de $X$los cuales no son preimagen de algun elemento de $Y$, es decir, elementos que no estan asociados a un elemento en $Y$por la funcion $g$, ie:     $$X_\infty = X \backslash \bigcup_{n=0}^\infty g^n(x)     $$donde $X_0 = X$. Tambien, definimos el complemento de este conjunto, $X_1$el cual contiene a todos los elementos que si estan asociados a un elemento de $Y$mediante la funcion $g$, ie:     $$X_1 = X \backslash X_\infty     $$Estos conjuntos son disjuntos entre si.
  \begin{figure}[H]
    \centering\includegraphics[width=30% ]{/books/1.jpg}\caption{Los colores en el grafico corresponden a los elementos relacionados segun la definicion de las funciones $f$y$g$. \newline En algunos libros dibulgativos se hace una analogia donde $X_\infty$son los elementos huerfanos de $X$debido a que son los elementos que no son pre-imagen de algun elemento en $Y$. De igual forma, se crean conjuntos de elementos padres, aquellos que si tienen una imagen en el conjunto codominio correspondiente a cada funcion asi como el conjunto de abuelos, etc.     }\end{figure}Ahora, construyamos una funcion $h: X \to Y$definida de la siguiente forma:     $$Si\;\;x \in X_\infty\;\Rightarrow h(x)=f(x)     $$$$Si\;\;x \in X_1\;\Rightarrow h(x)=g^{-1}(x)     $$y procedemos a probar que dicha funcion es una biyeccion entre $X$y $Y$. Para ello, notemos que dicha funcion es inyectiva por la definicion de $f$y $g$, puesto que cada elemento de $X_1$esta relacionado a un unico elementto $y \in Y$puesto que $f$es por deficion inyectiva. De manera analoga ocurre con $g$. Luego $h$es inyectiva. \newline Ahora, para probar que $h$es     \textbf{suprayectiva}: Una funncion $f:X\to Y$se dice que es suprayectiva si: \newline $\forall y \in Y;\;\exists x\in X \mid f(x)=y$, podemos notar que todos los elementos de $Y$son mapeados puesto que $\forall\; x \in X\; \exists ! y \in Y$tal que $f(x) = y$, por otro lado, por la definicion de $h$tenemos que todos los elementos de $f$tambien son mapeados por $h$, de igual forma ocurre con los elementos de $g$. Por lo que $h$es suprayectiva. \newline $\therefore \; h$es una biyeccion y tenemos entonces que $X \sim Y$
\item Dem: Sean $f:X \to Y$y $g:Y \to X$, ambas funciones inyectivas. Se tiene entonces que     $$f(X) = Y_1     $$donde $Y_1 \subset Y$. De esta forma $g$esta definida en $Y_1$, tal que     $$X_2 = g(Y_1) = g(f(X))     $$pero como $f$es inyectiva se tiene que $\forall x \in X, \exists ! y \in Y \mid f(x)=y$, esto a su vez significa que $X \sim Y_1$. Analogamente se puede usar un argumento similar tal que $Y_1 \sim X_2$por lo que por     \textbf{transitividad}: Es decir que si \newline $xRy \wedge yRz \Rightarrow xRz$$X \sim X_2$. De forma similar se puede verificar que $Y \sim Y_2 \subset Y$y continuando con este proceso podemos construir una susecion de subconjuntos     $$(X_n)_{n=0}^{\infty} \;\wedge\; (Y_n)_{n=0}^{\infty}     $$donde $X_0 = X$y $Y_0 = Y$. Esta sucesion de conjuntos se puede definir mediante     $$X_{n+1} = f(Y_n) \wedge Y_{n+1} = g(X_n)     $$lo que genera una cadena de conjuntos anidados $X_{n+1} \subset X_n$y $Y_{n+1} \subset Y_n$. \newline Recurriendo nuevamente a la inyectividad de $f$\begin{alert}{Esta demostracion no esta finalizada}\end{alert}
\end{itemize}La demostracion de este teorema es de gran relevancia ya que teniendo un criterio para determinar si dos conjuntos tienen la misma cardinalidad nos permite entonces establecer una relacion de  \textbf{orden total no estricto}: Una relacion de orden total no estricto es aquella relacion en la que se cuenta con la tricotomia y tambien con las propiedades de un orden parcial, es decir, es reflexiva, antisimetrica y finalmente, transitiva. entre los numeros cardinales, es decir, podemos comparar dichas cardinalidades, tambien se mostrara mas adelante, que $\delta(N) < \delta(R)$entre otros resultados importante
